% ==========================================================
% CUERPO PRINCIPAL DE LA TESIS
% Contenido alineado con el alcance/objetivos que definiste.
% ==========================================================

\chapter*{Resumen}
\addcontentsline{toc}{chapter}{Resumen}
% GUIA (tu proyecto):
% - Problema: movilidad con ceguera/baja visión, limitaciones herramientas tradicionales y sensores puntuales.
% - Objetivo de tesis: etapa 1 (captación + mapeo local compacto y estable en tiempo real).
% - Sensor: Intel RealSense D435i (RGB-D + IMU).
% - Output: matriz compacta baja resolución (ej 10x5) + tópicos ROS2 + interfaz serial.
% - Mapeo local con memoria temporal + compensación por orientación usando IMU.
% - Métricas: RMS error plano, latencia end-to-end, precisión/recall obstáculos, estabilidad temporal por celda, CPU/memoria.
% - Cierre: viabilidad como módulo perceptual para futura interfaz háptica (etapa 2, fuera de alcance).

\chapter{Introducción}

\section{Contexto y motivación}
% GUIA:
% - Describir problemática y motivación (como en tu texto).
% - Ventajas de profundidad vs sensores puntuales.
% - Idea general del proyecto en dos etapas y dónde cae esta tesis.

\section{Planteo del problema}
% GUIA:
% - Escribir el problema de ingeniería tal cual lo definiste:
%   sistema embebido + tiempo real + baja latencia + representación compacta estable.
% - Aclarar restricciones: portabilidad, cómputo limitado, estabilidad temporal.

\section{Objetivos}
\subsection{Objetivo del proyecto}
% GUIA:
% - Mencionar proyecto completo (2 etapas) solo para contexto.

\subsection{Objetivo de la tesis}
% GUIA:
% - Enunciar que esta tesis cubre "captación del entorno y mapeo local" únicamente.

\subsection{Objetivos técnicos}
% GUIA:
% - Pipeline RGB-D en ROS2 con D435i
% - Filtrado espacial/temporal
% - Agregación/reducción (matriz baja res)
% - Publicación ROS2 + salida serial
% - Mapa local con memoria + IMU para compensación orientación
% - Calibración/ajuste de parámetros por pruebas experimentales

\section{Alcance y exclusiones}
% GUIA:
% - IN: percepción + representación compacta + mapa local + IMU + publicación + serial.
% - OUT: diseño/implementación de interfaz háptica, actuadores, validación con usuarios finales, integración final portátil completa.
% - Plataforma: Raspberry Pi (modelo a determinar) o alternativa si no da tiempos.

\section{Organización del documento}
% GUIA:
% - Resumen de capítulos.

\chapter{Marco teórico}

\section{Percepción de profundidad por visión estéreo (RGB-D)}
% GUIA:
% - Conceptos de disparidad, profundidad, ruido, limitaciones típicas.
% - Qué entrega la D435i a nivel datos (sin entrar en datasheet si no querés).

\section{Representaciones espaciales locales para evitación de obstáculos}
% GUIA:
% - Mapas locales vs globales
% - Motivación de local/reactivo (mencionar línea Brock \& Khatib si corresponde a tu biblio).

\section{Filtrado espacial y temporal en mapas de profundidad}
% GUIA:
% - Qué se busca: reducción de ruido + estabilidad.
% - Principios generales (sin comprometerte a un filtro específico si todavía no lo cerraste).

\section{Compensación por orientación e integración IMU}
% GUIA:
% - Rol de la IMU: estimar orientación del sensor y compensar cambios de pose.
% - Conceptos generales: frames, rotaciones, alineación temporal.

\section{Métricas de evaluación para sistemas de percepción/mapeo}
% GUIA:
% - RMS en plano/geométrico, latencia end-to-end, precisión/recall, estabilidad temporal, consumo CPU/mem.

\chapter{Estado del arte}

\section{Dispositivos de asistencia con retroalimentación háptica}
% GUIA:
% - Los 3 que listaste (chaleco, bastón, pulsera).
% - Enfatizar limitaciones: resolución espacial, dependencia del brazo/cuerpo, alcance.

\section{Mapeo local y evitación de obstáculos en robótica móvil}
% GUIA:
% - Trabajos/ideas de uso de mapas locales, representaciones compactas.
% - Sensores RGB-D/estéreo con IMU y estabilidad ante movimiento.

\section{Discusión comparativa y posicionamiento}
% GUIA:
% - ¿Qué aporta tu tesis?: validación/adaptación de métodos robóticos al contexto de asistencia sensorial.
% - Por qué representación compacta (10x5 ej) es útil como interfaz hacia actuadores.

\chapter{Arquitectura del sistema}

\section{Arquitectura general}
% GUIA:
% - Diagrama de bloques:
%   D435i -> ROS2 (driver) -> pipeline percepción -> matriz compacta + mapa local -> tópicos -> serial -> (MCU etapa 2)
% - Aclarar que MCU/háptica es futura.

\section{Adquisición de datos con Intel RealSense D435i}
% GUIA:
% - Cómo se integra en ROS2 (paquete Intel).
% - Qué tópicos se usan (genérico si aún no fijaste nombres).

\section{Plataforma embebida}
% GUIA:
% - Raspberry Pi (modelo a determinar) como referencia.
% - Criterio de fallback a plataforma más potente si no cumple (ej Jetson Nano, como en tu plan).

\section{Diseño de interfaces}
\subsection{Interfaces ROS2}
% GUIA:
% - Tópicos publicados (matriz compacta, estado, debug, etc).
% - Logging/rosbags para reproducibilidad.

\subsection{Interfaz serial hacia microcontrolador}
% GUIA:
% - Formato general (matriz + timestamp/metadata si aplica).
% - Restricciones de ancho de banda y frecuencia.

\chapter{Diseño del método (pipeline de percepción)}

\section{Preprocesamiento de profundidad}
% GUIA:
% - ROI (si aplica), recorte por rango, manejo de valores inválidos.
% - Filtrado espacial (genérico si no definiste cuál).

\section{Agregación y reducción de dimensionalidad}
% GUIA:
% - Cómo pasás de mapa denso a matriz baja res (ej 10x5).
% - Qué representa cada celda (mínimo/mediana/promedio/etc) -> si no está decidido, dejar como "estrategias evaluadas".

\section{Representación compacta instantánea}
% GUIA:
% - Definición formal del output instantáneo por frame.

\section{Mapa local de corto alcance (memoria temporal)}
% GUIA:
% - Separación conceptual entre instantáneo y memoria.
% - Estrategia de fusión temporal y decaimiento/olvido.
% - Objetivo: robustez ante oclusiones y cambios de orientación.

\section{Integración de IMU para compensación}
% GUIA:
% - Cómo usás orientación para compensar/estabilizar el mapa local.
% - Alineación temporal IMU vs depth.

\section{Parámetros y calibración}
% GUIA:
% - Umbrales numéricos y criterios de decisión.
% - Cómo los vas a ajustar en etapa de calibración y pruebas.

\chapter{Implementación en ROS2}

\section{Estructura de paquetes y nodos}
% GUIA:
% - Paquete(s) implementados.
% - Nodos y responsabilidades (adquisición, procesamiento, serial, etc).
% - Configuración (params YAML) si aplica.

\section{Publicación y herramientas de experimentación}
% GUIA:
% - Uso de rosbag para registrar.
% - Trazabilidad de experimentos (versionado, configs).

\section{Implementación de la transmisión serial}
% GUIA:
% - Serialización de matriz.
% - Frecuencia, buffering, validación básica.

\chapter{Metodología experimental}

\section{Escenarios controlados}
% GUIA:
% - Referencias geométricas simples (plano) para error de profundidad.
% - Escenarios de obstáculos controlados (definir "obstáculo" operativamente).

\section{Métricas y criterios de evaluación}
% GUIA:
% (i) RMS en plano
% (ii) latencia end-to-end
% (iii) precisión/recall obstáculos
% (iv) estabilidad temporal por celda (variación)
% (v) consumo CPU/mem

\section{Instrumentación y procedimiento}
% GUIA:
% - Cómo medís latencia (timestamps, profiling).
% - Cómo medís CPU/mem (herramientas Linux).
% - Cómo rotulás ground truth en escenarios simples (sin inventar: describir metodología elegida).

\chapter{Resultados}

\section{Precisión de profundidad en escenarios controlados}
% GUIA:
% - Tablas/gráficos RMS por condición.

\section{Latencia end-to-end}
% GUIA:
% - Distribución/estadísticos y discusión vs requisitos.

\section{Detección de obstáculos}
% GUIA:
% - Precisión/recall, casos de falla.

\section{Estabilidad temporal del mapa local}
% GUIA:
% - Métrica por celda y ejemplos.

\section{Consumo de recursos}
% GUIA:
% - CPU/mem en plataforma objetivo.

\chapter{Discusión}
% GUIA:
% - Trade-offs: estabilidad vs latencia vs pérdida de detalle por reducción.
% - Limitaciones (hardware, ruido, movimiento).
% - Qué se demostró viable y qué queda abierto.

\chapter{Conclusiones}
% GUIA:
% - Resumen de aportes concretos:
%   (1) pipeline percepción en ROS2 con salida compacta
%   (2) mapa local con memoria + IMU
%   (3) evaluación cuantitativa con métricas
% - Conclusión de viabilidad para etapa 2.

\chapter{Trabajo futuro}
% GUIA:
% - Traducción a interfaz háptica (etapa 2)
% - Integración portátil completa
% - Validación con usuarios
% - Optimización/porting a hardware definitivo

% Fin del contenido principal
